% ============================================================
% 第 3 章 问题一：数据缺失条件下的感官评分
% ============================================================

\section{问题一：数据缺失条件下的感官评分}
\label{sec:q1}

\subsection{问题分析}

问题一要求给出 \texttt{3.xlsx} 中 Aroma 与 Flavour 两套、共约 40 个感官维度的 0--6 分评分。具体地，\texttt{3.xlsx} 的 Aroma 与 Flavour 各 15 维（共 30），\texttt{4.csv} 另含 20 个 taste/palate/overall 维度；本文对全部 50 维建模，并在结果中把 Aroma/Flavour 30 维作为核心、其余 20 维作为辅助输出。该任务有三个特殊之处。

第一，\textbf{没有可直接使用的等级标签}。\texttt{3.xlsx} 的 beer 评分卡是空白评分卡，\texttt{4.csv} 与 \texttt{7.csv} 提供的是标准化/中心化后的连续值，不是原始 0--6 分。若把标准化值直接四舍五入，会因样本集不同、中心化基准不同而产生不可比的分级。因此本文把任务拆成两层：先学习连续感官潜变量，再用有序阈值把它校准为 0--6 等级。

第二，\textbf{化合物侧存在缺失与高维共线}。\texttt{1.csv} 的化合物维度远大于样本量，且同源化合物高度相关；直接逐维回归会过拟合，缺失值处理不当还会把“缺失模式”误学成“风味信号”。因此模型必须显式使用缺失掩码，并在维度间共享信息。

第三，\textbf{感官维度之间存在相关结构}。\texttt{5.csv} 与 \texttt{6.csv} 表明不少感官维度共同变化；把 50 个维度完全独立建模会损失这部分结构。最终模型因此采用多输出/低秩共享结构，让不同感官维度共享化合物因子，并用集成融合补充非线性。

\subsection{评价指标}

设测试集共有 $n$ 个样本、$S$ 个感官维度。连续预测的评价指标为
\begin{align}
  \mathrm{MAE} &= \frac{1}{nS}\sum_{i=1}^{n}\sum_{s=1}^{S}\abs{\hat{y}_{i,s}-y_{i,s}},
  \label{eq:mae}\\
  \mathrm{RMSE} &= \sqrt{\frac{1}{nS}\sum_{i=1}^{n}\sum_{s=1}^{S}\left(\hat{y}_{i,s}-y_{i,s}\right)^2}.
  \label{eq:rmse}
\end{align}
等级预测采用二次加权 Kappa（QWK）与 Spearman 秩相关：
\begin{equation}
  \mathrm{QWK}=1-\frac{\sum_{k=0}^{6}\sum_{k'=0}^{6}W_{k,k'}O_{k,k'}}
  {\sum_{k=0}^{6}\sum_{k'=0}^{6}W_{k,k'}E_{k,k'}},
  \label{eq:qwk}
\end{equation}
其中 $O_{k,k'}$ 为真实等级 $k$ 被预测为 $k'$ 的观测频数，$E_{k,k'}$ 为对应期望频数，$W_{k,k'}=(k-k')^2/36$ 为二次权重\citep{cohen1968weighted}。Spearman 系数按
\begin{equation}
  \rho_s=1-\frac{6\sum_{i=1}^{n}\delta_i^2}{n(n^2-1)}
  \label{eq:spearman}
\end{equation}
计算，$\delta_i$ 为第 $i$ 个样本在两个排序中的名次差\citep{spearman1904proof}。

由于真实 0--6 标签不可得，本文在测试集上同时报告两类评价：其一是对连续监督信号的 MAE、RMSE 与 Spearman；其二是对“模拟等级”或专家复核等级的 QWK。前者衡量潜变量拟合，后者衡量分级校准；两者不可混为一谈。这也是本文把分级结果定位为“校准后的相对等级”而不是“绝对人工评分”的原因。

还需明确“参考分级”的含义：本文的参考分级是把连续监督信号按训练折内 $q_{0.01}/q_{0.99}$ 做稳健 min-max 后取整得到的离散结果，它是连续值的单调重标度，不含人工评分的绝对档位信息。因此后文报告的 QWK 衡量的是“同一连续信号在两种离散化下的一致性”，性质上接近离散化后的秩相关，不等于与人工评分逐档吻合。

\subsection{基线模型：逐维均值与线性回归}

最直接的基线是不使用化合物信息，只用训练集均值预测每个感官维度：
\begin{equation}
  \hat{y}_{i,s}^{(0)}=\bar{y}_s,\qquad
  \bar{y}_s=\frac{1}{N}\sum_{i=1}^{N}y_{i,s}.
  \label{eq:q1-mean}
\end{equation}
它给出“没有任何化合物信息”的下界。第二个基线是逐维岭回归，设 $\bbeta_s\in\R^{P}$ 为第 $s$ 个维度的回归系数，$c_s$ 为截距，$\bB=[\bbeta_1,\dots,\bbeta_S]$：
\begin{equation}
  \min_{\bB,\{c_s\}}
  \sum_{i=1}^{N}\sum_{s=1}^{S}\left(y_{i,s}-\vx_i^\top\bbeta_s-c_s\right)^2
  +\lambda\norm{\bB}_F^2.
  \label{eq:q1-ridge}
\end{equation}
岭回归通过共享特征矩阵与统一正则化降低小样本方差；PLS2 则把 $\bB$ 约束为低秩，用少量潜变量同时解释多个感官维度，是“多任务共享”的最简形式。二者都不使用缺失指示与风格信息，作为后续改进的对照\citep{hastie2009elements,tibshirani1996regression,wold2001pls}。

\subsection{改进模型：缺失感知的多输出回归与风格先验}

改进模型在岭回归/PLS2 的基础上加入两项数据事实。第一，化合物侧存在缺失：本文在训练折内做中位数插补，并把“某列是否缺失”作为 0/1 指示特征拼入 $\vx_i$，让模型区分“浓度低”与“未检测”。第二，\texttt{1.csv} 自带风格标签 $g(i)$，而同一化合物谱在不同风格下的感官期望并不相同；本文把风格先验拼入特征：
\begin{equation}
  \tilde{\vx}_i=\left[\vx_i;\ \bs_i\right],\qquad
  \bs_i\in\{0,1\}^{G}\ \text{或}\ \bs_i=\bar{\vy}_{g(i)},
  \label{eq:q1-style}
\end{equation}
其中 $\bar{\vy}_{g}$ 为训练折内风格 $g$ 的感官均值（target encoding），用留一或平滑避免泄漏。风格先验不是“用标签作弊”：它是 \texttt{1.csv} 的已有列，且在真实场景中风格在品评前已知；本文同时报告 chemical-only 与 chemical+style 的消融，量化它带来的提升。

\subsection{最终模型：集成融合与单调校准}

单一模型对共线化合物与非线性关系都不够稳健。最终模型先用五类基学习器得到 OOF 预测：PLS2（低秩多任务）、岭回归（线性）、ExtraTrees 与 RandomForest（非线性、多输出）、LightGBM（逐目标 boosting）。记第 $\nu$ 个基学习器对样本 $i$、维度 $s$ 的 OOF 预测为 $\hat{y}_{i,s}^{(\nu)}$，$\Mcal$ 为基学习器集合。本文比较两种融合：
\begin{equation}
  \hat{y}_{i,s}^{\mathrm{avg}}=\frac{1}{\abs{\Mcal}}\sum_{\nu\in\Mcal}\hat{y}_{i,s}^{(\nu)},
  \label{eq:q1-ensemble}
\end{equation}
\begin{equation}
  \hat{y}_{i,s}^{\mathrm{stack}}=\sum_{\nu\in\Mcal}\omega_{\nu,s}\,\hat{y}_{i,s}^{(\nu)}+b_s,
  \label{eq:q1-stack}
\end{equation}
其中 $\omega_{\nu,s}$ 在 OOF 预测上拟合。15 折结果显示等权融合的 $R^2=0.227$ 高于 Ridge 堆叠的 0.216，且折间方差更小，故最终采用式~\eqref{eq:q1-ensemble}；堆叠作为对照保留。融合的价值在于：线性模型与树模型在不同感官维度上各有胜负，等权平均让模型在不估计额外权重的情况下获得稳健性。

堆叠方案存在一处需要如实说明的乐观偏置：元学习器在训练折样本的折外预测上拟合，而这些折外预测来自训练时包含当前测试折样本的模型，故 0.216 这一数值偏高。即便带此偏置，堆叠仍未超过等权融合的 0.227，这进一步支持选择等权融合作为最终方案。

连续预测存在向均值收缩的问题，直接做 min-max 会把 0--6 分级挤到中间几档。本文在训练折内拟合单调校准 $g_s$（isotonic regression，保序回归），把 OOF 预测映射到真实连续值，再做稳健 min-max：
\begin{equation}
  \hat{k}_{i,s}=\operatorname{round}\left(
  6\,\frac{g_s(\hat{y}_{i,s})-q_{0.01,s}}{q_{0.99,s}-q_{0.01,s}}
  \right),\qquad \hat{k}_{i,s}\in\{0,\dots,6\},
  \label{eq:q1-grade}
\end{equation}
其中 $q_{0.01,s},q_{0.99,s}$ 为训练折内 $g_s(\hat{y}_{\cdot,s})$ 的 1\%/99\% 分位数，越界值截断。保序回归只要求单调，不假设线性，能纠正收缩而不改变样本排序；min-max 只负责把校准后的连续值锚定到 0--6。quantile 与 Gaussian latent 两种校准作为敏感性分析，用于检查分级结论是否依赖校准假设。

\figauto{fig_q1_framework.pdf}{问题一的分级评分框架：基线、改进与最终模型}{fig:q1framework}

图~\ref{fig:q1framework} 概括了三种路线的输入、输出与衔接关系：基线只用均值/线性回归；改进模型加入缺失指示与风格先验；最终模型用集成融合与保序校准。三者在同一数据划分上比较，避免“换了数据又换了模型”的不可比问题。

\subsection{实验设置}

本文按样本分层划分训练、验证与测试集，所有随机过程固定种子；缺失插补、标准化、保序校准与超参数选择只在训练/验证集内完成。为检验缺失条件下的稳健性，在测试集上按不同比例随机遮蔽化合物，并额外模拟分位数驱动的非随机缺失（MAR），比较三种路线的退化曲线。主要超参数包括正则化强度 $\lambda$、PLS 潜变量数、树模型叶子数/树数、LightGBM 学习率、集成权重与校准分位数；其搜索范围与最终取值记录在附录表~\ref{tab:appendix-config} 中。本文使用 5 折 $\times$ 3 次重复（种子 42），MCAR 5/10/20/30\%，MAR 由 \texttt{ethanol..v.v.} 的分位数驱动。所有超参数均在外层训练折内用内层交叉验证选择、不跨折固定：岭回归 $\alpha\in[10^{0},10^{6}]$（13 点对数网格、内层 5 折）、PLS 成分数 $2$--$20$（内层 5 折）、ExtraTrees 与 RandomForest 叶节点 $\in\{1,2,4\}$（内层 3 折），树模型最终 500 棵；LightGBM 固定 400 棵、学习率 0.03、叶子 15。堆叠元学习器为 RidgeCV（$\alpha\in[10^{-3},10^{2}]$、内层 3 折）。各折实际取值记录在 \texttt{data/q1/model\_summary.json} 的 \texttt{hyperparameters} 字段。

对照设置包括：（1）逐维均值基线；（2）逐维岭回归；（3）式~\eqref{eq:q1-ridge} 的多输出岭回归/PLS2；（4）式~\eqref{eq:q1-ensemble} 的等权融合；（5）去掉风格先验、去掉保序校准、或去掉树模型的消融版本。通过这一组对照，可以分辨“多任务共享”“风格先验”和“保序校准”各自的贡献，而不是只报一个最终分数。

\subsection{结果与分析}

表~\ref{tab:q1-results} 给出三种路线在重复交叉验证上的连续拟合与分级一致性指标。最终模型相对基线 Ridge 的 $R^2$ 提升 0.065（0.162$\to$0.227），Spearman 提升 0.059（0.373$\to$0.432）。

\begin{table}[htbp]
  \centering
  \caption{问题一三种路线的交叉验证结果}
  \label{tab:q1-results}
  \small
  \begin{tabular}{@{}l c c c c c@{}}
    \toprule
    模型 & $R^2$ & MAE & RMSE & Spearman $\rho_s$ & QWK \\
    \midrule
    基线：逐维岭回归 & 0.1617 & 0.3622 & 0.4946 & 0.3728 & — \\
    改进：岭回归 + 风格先验 & 0.1855 & 0.3542 & 0.4855 & 0.3927 & — \\
    最终：五模型等权融合 + 保序校准 & 0.2270 & 0.3415 & 0.4695 & 0.4318 & 0.3756 \\
    \bottomrule
  \end{tabular}
\end{table}

表中 R²/RMSE/MAE/Spearman 为 5 折 $\times$ 3 次重复的折间均值；QWK 只对最终分级校准报告，基线与改进未做分级校准。最终模型的折间标准差为 R² 0.021、Spearman 0.024、RMSE 0.018、MAE 0.011；基线与改进的标准差见支撑材料。风格先验在岭回归上稳定提升 R² $+0.024$（0.1617$\to$0.1855），在融合模型上提升 $+0.012$（0.2270$\to$0.2388）；由于融合模型已通过化学特征间接吸收风格，最终交付仍以 chemical-only 融合为主，style 版本作为消融报告。

从指标结构看，连续指标与分级指标回答的是不同问题：MAE/RMSE 衡量连续潜变量是否贴合监督信号，QWK 衡量分级是否与等级顺序一致。最终模型的 14/50 个目标 $R^2>0.3$，其中 acidity（0.690）、alcohol（0.678）、lactic（0.648）、acetic（0.635）、bitternes（0.604）最强；barnyard（$-0.039$）、stale\_hops（$-0.032$）、clove（$-0.005$）最弱。复杂/微生物/阈值型维度基本不可由 231 个化学特征预测，这与化学合理性文献一致，也解释了为什么 Q2 的纯数据 Top-3 会在这些维度上选到代理。

\figauto{q1_target_r2_heatmap.pdf}{逐感官维度的交叉验证 $R^2$}{fig:q1target}

在缺失稳健性方面，自然缺失（73/57,750 $=0.126\%$）下最佳模型的 $R^2$ 为 0.2264；MCAR 5\%/10\%/20\%/30\% 时 $R^2$ 相对下降 3.46\%/6.33\%/11.46\%/19.33\%，RMSE 相对上升 0.70\%/1.15\%/2.28\%/3.75\%；MAR（24.77\%）的 $R^2$ 下降 10.89\%，与同量级 MCAR 接近。30\% MCAR 内退化可控，说明缺失指示与折内中位数插补有效；但没有证据表明模型能完全消除非随机缺失偏差，这一限制在讨论中保留。

\figauto{q1_missingness.pdf}{缺失率—性能曲线（MCAR 与 MAR）}{fig:q1missing}

\subsection{讨论}

问题一的主要困难不是模型复杂度，而是标签不可直接观测。本文用“连续监督 + 有序校准”绕开了空白评分卡的限制，但这也意味着最终 0--6 分数的绝对水平依赖校准锚点，不能声称与人工评分逐档完全一致。更稳妥的表述是：在给定校准方案下，模型能保持样本间的相对顺序与等级单调性；若要报告绝对等级精度，需要原始评分或专家复核。

其次，低秩结构假定不同感官维度由少量共享因子驱动，这与 \texttt{5.csv}、\texttt{6.csv} 的相关结构一致，但若某些维度由独立机制主导（例如特定缺陷阈值），低秩约束可能欠拟合。解决办法是允许部分维度使用独立残差项，或在低秩模型与逐维模型之间做模型平均；这部分作为后续改进方向。

最后，缺失机制未知。本文的掩码损失在随机缺失下是无偏的，在非随机缺失下只能减轻偏差，不能消除。若后续能获得缺失原因的元数据（例如某化合物未检测而非浓度为零），应把缺失原因作为额外变量纳入模型。
